\documentclass[12pt,a4paper]{article}

\usepackage[utf8]{inputenc}
\usepackage[T1]{fontenc}
\usepackage[french]{babel}
\usepackage{lmodern}
\usepackage[a4paper,margin=2.5cm]{geometry}
\usepackage{hyperref}
\usepackage{xcolor}
\usepackage{array}
\usepackage{booktabs}
\usepackage{makecell}
\usepackage{tabularx}
\usepackage{enumitem}
\usepackage{amsmath,amssymb}
\usepackage{tcolorbox}
\usepackage{float}
\usepackage{parskip}

\definecolor{accent}{HTML}{2A5A8A}
\definecolor{warn}{HTML}{B85C00}
\definecolor{thbg}{HTML}{EEF3F8}
\definecolor{okbg}{HTML}{EEF9EE}
\definecolor{ok}{HTML}{2A7A2A}

\hypersetup{colorlinks=true, linkcolor=accent, urlcolor=accent, citecolor=accent}

\tcbuselibrary{skins,breakable}
\newtcolorbox{insight}{
  colback=thbg, colframe=accent, boxrule=0.8pt,
  left=8pt, right=8pt, top=6pt, bottom=6pt
}
\newtcolorbox{attention}{
  colback=okbg, colframe=ok, boxrule=0.8pt,
  left=8pt, right=8pt, top=6pt, bottom=6pt
}
\newtcolorbox{limite}{
  colback=yellow!8, colframe=warn, boxrule=0.8pt,
  left=8pt, right=8pt, top=6pt, bottom=6pt
}

\renewcommand{\arraystretch}{1.5}

\title{\textbf{Résultats et analyse du benchmark RTS-LLM}}
\author{Samir El Khoumri, Aly Hachem-Reda, Sami Chbicheb}
\date{Mai 2026}

\begin{document}

\maketitle
\tableofcontents
\newpage

% ============================================================
\section{Introduction}
% ============================================================

Nous avons vu que la vérité terrain est l'élément central de toute évaluation
d'un outil de sélection de tests : sans elle, il est impossible de savoir si
l'outil sélectionne les bons scénarios ou s'il se trompe. Elle constitue la
référence objective à laquelle toute sélection automatique doit être comparée.

Pour exploiter cette vérité terrain, il faut un \textbf{benchmark} : un
protocole d'évaluation structuré qui soumet l'outil à un ensemble de situations
réelles et mesure ses performances de manière reproductible. Le benchmark joue
ici un rôle clé. Il ne s'agit pas simplement de vérifier que l'outil
``fonctionne'' dans quelques cas favorables, mais de le confronter à la
diversité des commits réels --- des modifications triviales aux refactorisations
complexes --- pour obtenir une image honnête de ce qu'il sait faire et de ce
qu'il rate. C'est cette image qui permet de prendre des décisions éclairées sur
l'utilisation de l'outil en production et sur les axes d'amélioration prioritaires.

Dans notre cas, le benchmark sert également à répondre à une question précise :
l'ajout d'une analyse statique du code apporte-t-il réellement quelque chose par
rapport au LLM utilisé seul ? Les résultats qui suivent quantifient cet apport
commit par commit, sur les trois métriques standard : précision, rappel et F1.

% ============================================================
\section{Résultats quantitatifs}
% ============================================================

Le benchmark compare deux configurations de RTS-LLM sur 38 commits du projet
\textbf{spring-cucumber-rest-api}, en utilisant le modèle Mistral~7B-instruct
(q3\_K\_M) via Ollama sur GPU local (NVIDIA RTX A500, 4\,Go VRAM).

\begin{itemize}[noitemsep]
  \item \textbf{LLM seul} : le diff et tous les scénarios BDD sont soumis
        directement au modèle, qui sélectionne seul les scénarios à retester.
  \item \textbf{Hybride} : un outil d'analyse de code identifie d'abord
        quels scénarios sont reliés aux méthodes modifiées en remontant
        les chaînes d'appel ; le LLM reçoit uniquement ces candidats
        et décide lesquels sont réellement impactés.
\end{itemize}

\subsection{Métriques moyennes}

\begin{table}[H]
\centering
\begin{tabular}{lcccc}
\toprule
\textbf{Configuration}
  & \makecell{\textbf{Commits}\\\textbf{gt\,>\,0}}
  & \makecell{\textbf{Précision}\\\small(commits actifs)}
  & \makecell{\textbf{Rappel}\\\small(tous commits)}
  & \makecell{\textbf{F1 moyen}\\\small(tous commits)} \\
\midrule
LLM seul                 & 24 & 0{,}487 & 0{,}368 & 0{,}227 \\
Hybride (statique + LLM) & 17 & \textbf{0{,}789} & 0{,}372 & \textbf{0{,}378} \\
\midrule
Gain                     & --- & $\mathbf{+0{,}302}$ & $+0{,}004$ & $\mathbf{+0{,}151}$ \\
\bottomrule
\end{tabular}
\caption{Métriques moyennes sur les commits avec \texttt{gt\_size > 0}
         (vérité terrain non vide). Modèle : Mistral 7B-instruct q3\_K\_M.}
\label{tab:resultats}
\end{table}

\subsection{Distribution du rappel par commit}

\begin{table}[H]
\centering
\begin{tabular}{lcc}
\toprule
\textbf{Tranche de rappel}
  & \textbf{LLM seul} ($n = 24$)
  & \textbf{Hybride} ($n = 17$) \\
\midrule
$R = 0$ \quad (aucun scénario impacté détecté)  & 13~commits (54\,\%) & 6~commits (35\,\%) \\
$0 < R < 0{,}5$ \quad (détection partielle)     &  2~commits ( 8\,\%) & 5~commits (29\,\%) \\
$0{,}5 \leq R < 1$ \quad (bonne détection)      &  2~commits ( 8\,\%) & 1~commit  ( 6\,\%) \\
$R = 1$ \quad (détection complète)              &  7~commits (29\,\%) & 5~commits (29\,\%) \\
\bottomrule
\end{tabular}
\caption{Distribution du rappel sur les commits avec \texttt{gt\_size > 0}.}
\label{tab:distribution}
\end{table}

% ============================================================
\section{Interprétation des résultats}
% ============================================================

\subsection{Pourquoi le mode hybride obtient une précision bien meilleure}

Le gain en précision est le résultat le plus marquant du benchmark
($+0{,}302$, soit $+62\,\%$ relatif). Il s'explique directement par le rôle
de l'analyse statique comme pré-filtre.

\textbf{En mode LLM seul}, le modèle reçoit la totalité des scénarios BDD
sans savoir lesquels exécutent réellement le code modifié. Il se base alors
sur la \emph{ressemblance de mots} entre le diff et les descriptions de
scénarios : il sélectionne un scénario si sa description contient des termes
proches du diff (noms de méthodes, type d'opération, champs modifiés).
Ce raisonnement produit des \textbf{faux positifs} : des scénarios dont
le texte ressemble au diff mais qui n'exécutent pas le code modifié.
La précision de 0{,}487 signifie qu'en moyenne, sur deux scénarios sélectionnés
par le LLM seul, un seul est réellement impacté.

\textbf{En mode hybride}, seuls les scénarios dont le code de test appelle
(directement ou indirectement) une méthode modifiée sont transmis au LLM.
Ce filtre préalable élimine les faux positifs par ressemblance de mots :
un scénario ne peut figurer dans la liste que s'il est prouvé qu'il passe par
le code modifié. Le LLM reçoit en plus le chemin d'appel reliant chaque scénario
à la méthode modifiée, ce qui lui permet de confirmer ou d'écarter chaque cas
avec des informations concrètes plutôt que de deviner. La précision de 0{,}789
signifie que près de quatre scénarios sélectionnés sur cinq sont réellement impactés.

\begin{insight}
\textbf{En résumé :} l'analyse de code sert de filtre objectif qui supprime
les faux positifs que le LLM produirait par ressemblance de mots.
Le LLM, muni du chemin d'appel concret, peut ensuite juger si chaque
scénario est vraiment impacté par le changement.
\end{insight}

\subsection{Pourquoi le rappel n'augmente presque pas}

Malgré la nette amélioration de précision, le rappel reste quasi identique
(0{,}368 \emph{vs} 0{,}372, $+0{,}004$). Ce résultat peut sembler paradoxal :
si l'analyse statique « trouve » les bons scénarios, pourquoi ne pas mieux
les sélectionner ?

\textbf{Le LLM introduit ses propres faux négatifs.} En mode hybride,
si le modèle juge qu'un lien est trop indirect ou que la modification est
mineure, il écarte ce scénario de la sélection finale — même si l'analyse
de code l'avait correctement identifié comme impacté. Le gain apporté par
le filtre de code est donc partiellement annulé.

\textbf{L'analyse de code produit elle aussi des faux négatifs.} Les méthodes
ajoutées par le diff, les appels faits par réflexion ou par héritage, et
les classes supprimées ne sont pas détectés par l'outil d'analyse.
Si un scénario est impacté par l'un de ces cas, il ne figurera jamais dans
la liste des candidats et sera manqué, quelle que soit la décision du LLM.

\textbf{La baisse des sélections vides compense ces faux négatifs.}
Le seul gain de rappel visible est la baisse du taux de sélection vide
(de 54\,\% à 35\,\%). Quand le LLM seul ne voit pas de lien évident, il
ne sélectionne rien ; quand le mode hybride lui fournit un chemin d'appel
concret, il sélectionne au moins quelques candidats. C'est cet effet qui
explique les $+0{,}004$ de rappel moyen : les scénarios récupérés grâce au
filtre compensent tout juste les faux négatifs introduits par le LLM.

\subsection{Ce que révèle la distribution du rappel}

La distribution (Tableau~\ref{tab:distribution}) donne une image plus fine
que les seules moyennes.

\textbf{Le mode LLM seul fonctionne en tout ou rien.} 54\,\% des commits
produisent $R = 0$ et 29\,\% produisent $R = 1$, avec très peu de cas
intermédiaires (16\,\%). Soit le LLM ne voit aucun lien et ne sélectionne
rien, soit il identifie un lien évident et sélectionne correctement tous
les scénarios impactés. Les détections partielles sont rares.

\textbf{Le mode hybride couvre davantage de cas intermédiaires.}
La proportion de commits avec $0 < R < 1$ passe de 16\,\% à 35\,\%.
Le filtre de code fournit une liste de candidats que le LLM affine :
cette sélection partielle vaut mieux que rien — même un rappel de 0{,}3
ou 0{,}5 réduit déjà fortement le nombre de tests à exécuter.

\textbf{Le taux de détection complète ($R = 1$) est identique (29\,\%).}
Les commits sur lesquels l'outil trouve tous les scénarios impactés sont
aussi nombreux dans les deux modes. Ce sont des commits avec un diff court
ciblant une ou deux méthodes : le lien avec les scénarios est si évident
que le LLM le détecte directement, sans avoir besoin du filtre de code.

\begin{attention}
Les 7 commits avec $R = 1$ en mode LLM seul sont :
\texttt{34eb988}, \texttt{53f6cc1}, \texttt{5799273}, \texttt{86a39d5},
\texttt{afb6fc0}, \texttt{bd99a7c}, \texttt{fcd81b4}.
Ils partagent tous un diff court ciblant une méthode dont le nom ou le rôle
est suffisamment parlant (\texttt{ControllerAdvice}, \texttt{authenticate},
\texttt{getAll\ldots}) pour que le LLM fasse le lien sans aide extérieure.
\end{attention}

\subsection{Lecture synthétique : deux profils de commits}

L'ensemble des résultats converge vers un constat : les performances de
l'outil dépendent fortement du \textbf{profil du commit}, indépendamment
du mode utilisé.

\paragraph{Commits avec un diff court et ciblé.}
Le diff modifie une ou deux méthodes dont le rôle est précis et limité.
Peu de scénarios passent par ce code, donc peu sont impactés. Le lien
entre la modification et les scénarios concernés est clair, et les deux
modes trouvent correctement tous les scénarios ($R = 1$). C'est le cas
idéal pour RTS-LLM.

\paragraph{Commits avec un diff long touchant du code partagé par tous les scénarios.}
Le diff modifie du code exécuté par presque tous les scénarios
(gestion des requêtes HTTP, authentification, gestionnaire d'erreurs global).
Presque tous les scénarios sont donc impactés, mais l'outil n'ose pas
tous les sélectionner. Le résultat est souvent une sélection vide
(LLM seul) ou partielle (hybride), avec un rappel très faible.

\begin{limite}
Quand la modification touche du code partagé par tous les scénarios,
l'outil ne peut pas remplacer l'exécution complète de la suite de tests.
Il reste cependant utile pour la majorité des commits du quotidien,
dont le diff est court et ciblé sur une fonctionnalité précise.
\end{limite}

% ============================================================
\section{Perspectives d'amélioration}
% ============================================================

Les résultats identifient deux axes de progrès prioritaires.

\textbf{Réduire le taux de sélection vide ($R = 0$).} C'est le principal
levier pour améliorer le rappel. Trois pistes sont envisageables :
(1)~utiliser un modèle plus grand (13B, 34B) qui serait moins réticent à
sélectionner quand le diff est long ; (2)~augmenter la limite de taille du
diff transmis (actuellement 4\,000 caractères) pour que le LLM dispose de
plus de contexte ; (3)~reformuler le prompt du mode hybride pour que le LLM
ne génère pas de réponses hors-sujet quand le contexte est partiel.

\textbf{Améliorer la précision du mode LLM seul.} Une précision de 0{,}487
signifie qu'un scénario sélectionné sur deux est un faux positif. Fournir
au LLM quelques informations sur la structure du code — même sans analyse
complète des chaînes d'appel — pourrait limiter les sélections basées
uniquement sur la ressemblance de mots.

\end{document}
